\section{CUDA GPU并行编程：向量归约中的线程利用率优化}

\subsection{引言}
在CUDA应用程序开发中，提升性能的核心策略之一是\textbf{增加每个线程的工作量}（Work perThread）。
这一优化不仅适用于向量加法，对向量归约（VectorReduction）同样至关重要。在基础的向量归约实现中，往往存在
严重的线程利用不足问题：在每一轮迭代中，仅有一半的线程处于活跃状态，另一半则处于空闲等待状态。这种资源浪费直接导致了流式
多处理器（SM）的占用率（Occupancy）低下。本节将详细讲解如何通过让每个线程处理更多元素来消除空闲线程，从而显著
提升硬件利用率与整体性能。

\subsection{问题分析：线程利用不足}
在之前的基准实现中，我们设定块大小（BlockSize）为256个线程，每个块处理256个元素。此时，线程数量严格等于元
素数量。然而，在归约操作的第一轮迭代中，只有偶数ID的线程执行加法操作，奇数ID的线程完全空闲。这意味着：
\begin{itemize}
    \item 50\%的计算资源被浪费；
    \item SM的实际占用率远低于理论峰值；
    \item 基准测试显示，包含分支发散（Warp Divergence）的内核执行时间为283微秒，即使移除条件分支优化后仍需176微秒，性能仍有巨大提升空间。
\end{itemize}

\subsection{优化策略：增加每线程工作量}
\subsection{核心思想}
我们将每个块处理的元素数量从256增加到512，但保持块大小仍为256个线程。这意味着\textbf{每个线程负责处理2个元素}。通过这种方式，所有256个线程在第一阶段都能满载工作，彻底消除了空闲线程。

\subsection{两阶段执行流程}
新的内核函数执行分为两个逻辑步骤：
\begin{enumerate}
    \item \textbf{预归约（Pre-reduction）}：每个线程将分配给它的第二个元素（偏移量为 blockDim）加到第一个元素上。例如，线程 $i$ 执行 $data[i] += data[i + blockDim]$。此步骤将512个元素归约为256个有效元素，且所有线程均参与计算。
    \item \textbf{标准归约}：经过第一步处理后，数据规模变回256个元素，这与基准版本的输入完全一致。因此，后续可直接复用原有的归约逻辑，无需修改核心循环代码。
\end{enumerate}

\subsection{网格与索引调整}
由于每个块处理的元素翻倍，启动配置需相应调整：
\begin{itemize}
    \item \textbf{全局索引计算}：$index = (blockIdx.x \times blockDim.x + threadIdx.x) \times 2$。乘以2是因为每个线程现在跨越两个元素位置。
    \item \textbf{网格大小（Grid Size）}：变为原来的一半。因为总元素数不变，但每个块的处理能力翻倍，所需的块数量自然减半。这减少了SM调度的波次（Waves），有助于减少尾效应并提升整体吞吐。
\end{itemize}

\subsection{性能分析与实测结果}
\subsection{执行时间对比}
通过Nsight Compute等工具分析，优化效果显著：
\begin{itemize}
    \item 基准版本（含分支）：283 $\mu s$
    \item 移除分支版本：176 $\mu s$
    \item \textbf{增加线程工作量版本：100 $\mu s$}
\end{itemize}
相比基准版本，性能提升接近3倍；相比仅移除分支的版本，仍有约43\%的性能提升。

\subsection{占用率与内存行为}
\begin{itemize}
    \item \textbf{占用率}：提升至接近90\%，证明SM资源得到了充分利用。
    \item \textbf{缓存流量}：L1到L2缓存的数据写入量约为19MB。虽然计算效率提升，但全局内存访问模式仍有优化空间，这为后续引入共享内存（Shared Memory）优化埋下伏笔。
\end{itemize}

\subsection{进阶优化探讨}
\subsection{循环遍历顺序的影响}
实验表明，将归约循环从"大步长递减"（Stride = BlockDim/2 $\to$1）改为"小步长递增"（Stride = 1 $\to$ BlockDim/2）在本案例中反而导致性能下降（从100$\mu s$ 劣化至150 $\mu s$）。
尽管某些文献建议小步长优先有利于缓存局部性，但在未使用共享内存且数据跨度较大的场景下，大步长递减可能更符合当前的内存合并访问模式。\textbf{注意：最优策略高度依赖于具体硬件架构与数据规模，务必实测验证。}

\subsection{位运算替代除法}
在归约循环中，频繁使用除法和取模运算代价高昂。利用位运算可进一步微调性能：
\begin{itemize}
    \item 右移一位（\verb|>> 1|）等价于除以2；
    \item 左移一位（\verb|<< 1|）等价于乘以2。
\end{itemize}
实测显示，将步长更新逻辑替换为位运算后，获得了约2\%的额外性能提升。虽然幅度不大，但在大规模并行计算中，消除百万级线程
的整数除法开销仍是良好的编程实践。

\subsection{总结与展望}
本节证明了“增加每线程工作量”是解决向量归约中线程闲置问题的有效手段。通过将每块处理元素翻倍，我们实现了全线程活跃，大幅
降低了执行时间并提升了SM占用率。然而，当前实现仍受限于全局内存带宽。下一阶段，我们将结合\textbf{共享内存}技术，
减少对L2缓存及全局内存的访问压力，进一步挖掘GPU性能潜力。建议学习者在掌握本节内容后，尝试将共享内存与线程工作量优
化相结合，以获得最佳实践效果。

\section*{附录：关键代码片段}

\begin{lstlisting}[caption={优化后的向量归约内核（增加每线程工作量）}]
__global__ void reduceOptimized(float *g_idata, float *g_odata, unsigned int n) {
    extern __shared__ float sdata[]; // 预留共享内存，本节暂未启用
    
    // 每个线程处理2个元素，全局索引需乘2
    unsigned int tid = threadIdx.x;
    unsigned int i = (blockIdx.x * blockDim.x + threadIdx.x) * 2;
    
    // 第一阶段：预归约，将所有线程激活
    // 将后半部分数据加到前半部分对应位置
    if (i + blockDim.x < n) {
        g_idata[i] += g_idata[i + blockDim.x];
    }
    
    // 第二阶段：标准归约（此时有效数据量已减半）
    // 后续逻辑与原256元素归约完全相同
    for (unsigned int stride = blockDim.x / 2; stride > 0; stride >>= 1) {
        __syncthreads();
        if (tid < stride) {
            g_idata[i] += g_idata[i + stride];
        }
    
    if (tid == 0) {
        g_odata[blockIdx.x] = g_idata[i];
    }
\end{lstlisting}

